% Chapter 4 -- Methodology I: Detector, Crop Pipeline, and the Conceded Ordinal Engine
% Self-contained section file: no preamble, no \documentclass.
% Mirrors the template's Chapter 3 (Sec 3.2 Data Preprocessing + Sec 3.3 Modeling).
% HONESTY: no experimental results are reported here. This chapter specifies the
% engineering of the engine and its expected/gate-conditioned behaviour only.

\chapter{Methodology~I: Detector, Crop Pipeline, and the Conceded Ordinal Engine}
\label{ch:methodology-engine}

\section{Scope and a standing concession}
\label{sec:eng-scope}

This chapter specifies the perception \emph{engine} of our system: the binary
pain detector (Phase~A), the face-crop and eye-alignment preprocessing that
feeds it forward, the frozen DINOv2 backbone, the five per-action-unit (per-AU)
ordinal heads, the distributional decode, and the training procedure that fits
the heads. It mirrors the modelling and preprocessing chapter of the reference
template, but it does so under a concession that governs every subsection below
and that we restate without euphemism: \emph{the engine is plumbing, not a
contribution}. The frozen-backbone-plus-ordinal-heads design is a standard,
defensible low-data transfer-learning recipe; we adopt it because it is correct
for our regime, not because it is novel, and we never claim otherwise. The
contribution of this work lives in a portable method built on top of the
engine's outputs: the power-aware, per-AU capture-condition
confound-attribution protocol, which is the headline and is developed in a later
chapter. Supporting it are a guarded VLM-as-AU-rater $\kappa$ reliability check
(inspected-not-validated, with its result pending an independent vet anchor) and
the cited trustworthiness wrapper, which are developed in later chapters.
Everything specified here exists only to
produce calibrated scores for those methods to interrogate.

Two framing commitments follow directly from that concession and recur
throughout the chapter. First, the \emph{validated spine of v1 is binary}: the
per-image pain decision (Section~\ref{sec:eng-detector}) is the only output of
this chapter that is later subjected to a validated sensitivity/specificity
claim, and it is validated on vet-confirmed labels only. Second, the
\emph{graded $0$--$10$ layer ships inspected-not-validated}: the five ordinal
heads, their summed score, and the distributional pmf over that sum
(Sections~\ref{sec:eng-corn}--\ref{sec:eng-decode}) are built, inspected, and
reported qualitatively, but they never enter a validated-claim sentence. The
word ``graded'' is struck from every such sentence by construction, and the
project's headline must hold with the graded layer dropped. We flag the
engineering hooks for this discipline as we reach them.

No numerical results appear in this chapter. Where we state a quantity --- an
expansion margin, an alignment tolerance, an operating point --- it is a
\emph{design value} or a \emph{gate threshold} to be confirmed against
pilot output before freezing, not a measured outcome. Literature figures
(e.g.\ the Feline Grimace Scale validation anchors) are attributed to their
sources and used only to motivate design choices.

Figure~\ref{fig:eng-pipeline} gives the end-to-end view that the remaining
sections fill in.

\begin{figure}[t]
\centering
\begin{tikzpicture}[
  node distance=6mm and 9mm,
  box/.style={draw, rounded corners, align=center, minimum height=9mm,
              inner sep=3pt, font=\footnotesize, text width=20mm},
  conc/.style={box, fill=black!6},      % conceded plumbing
  val/.style={box, draw=black, very thick},
  arr/.style={-{Latex[length=2mm]}, thick},
  lbl/.style={font=\scriptsize\itshape, text=black!55}
]
\node[box] (img) {Frame\\(image)};
\node[conc, right=of img] (det) {RF-DETR\\pain box\\(Phase A)};
\node[conc, right=of det] (crop) {margin-expand\\+ eye-align\\crop $518^2$};
\node[conc, right=of crop] (bb) {frozen\\DINOv2\\ViT-S/14\\\textbf{(cached)}};
\node[conc, below=12mm of bb] (heads) {5 per-AU\\CORN heads};
\node[conc, left=of heads] (pmf) {per-AU\\pmf $\{0,1,2\}$};
\node[conc, left=of pmf] (conv) {convolve\\$\rightarrow$ 0--10\\sum pmf};
\node[val, left=of conv] (dec) {$\geq 0.39$\\point\\decision};

\draw[arr] (img) -- (det);
\draw[arr] (det) -- (crop);
\draw[arr] (crop) -- (bb);
\draw[arr] (bb) -- (heads);
\draw[arr] (heads) -- (pmf);
\draw[arr] (pmf) -- (conv);
\draw[arr] (conv) -- (dec);

\node[lbl, below=1mm of det.south] {Gate 1/3 folds};
\node[lbl, below=1mm of crop.south] {Gate 5 (NME)};
\node[lbl, above=1mm of heads.north] {Gate 4 (decode)};
\node[lbl, below=1mm of dec.south, text width=24mm, align=center]
   {\emph{inspected-not-validated}};
\end{tikzpicture}
\caption[End-to-end engine pipeline]{End-to-end engine pipeline. A frame is
localised by the Phase~A detector, the box is margin-expanded and (if Gate~5
requires) eye-aligned into a fixed $518\times518$ crop, encoded once by the
frozen DINOv2 backbone (features cached to disk), and read by five per-AU CORN
heads. Soft per-AU probability mass functions are convolved into a distribution
over the $0$--$10$ sum; a hard decode at the $0.39$ ratio yields the point
decision. Shaded blocks are \emph{conceded plumbing}; the heavy-bordered
decision block and the binary pain decision in Section~\ref{sec:eng-detector}
are the only validated outputs. The $0$--$10$ sum and its pmf ship
inspected-not-validated. Letters in italics mark the gate that governs each
stage.}
\label{fig:eng-pipeline}
\end{figure}

\section{Phase A: the binary pain detector}
\label{sec:eng-detector}

Phase~A is the binary spine of the system and the only stage that produces a
validated number. Its job is twofold: to emit a per-image \emph{pain /
no-pain} decision, and to supply a tight head box that becomes a
\emph{candidate} face crop for the rest of the engine. The decision is the
contribution-bearing output; the crop hand-off is conditional on the alignment
audit of Section~\ref{sec:eng-crop} and is treated as a risk until that audit
passes.

\subsection{Architecture choice}
\label{sec:eng-detector-arch}

We adopt an object-detection formulation rather than a crop-then-classify
pipeline, because the detector doubles as the downstream face cropper and a
separate localiser would be wasted work. The primary architecture is
\textbf{RF-DETR-Nano}, stepping up to RF-DETR-Small only if pain recall
plateaus below target; the runner-up, used as a fallback if RF-DETR training is
unstable on this small, imbalanced corpus, is \textbf{YOLOv11s}. Both are
fine-tuned from COCO-pretrained weights --- the COCO ``cat'' category makes
transfer expected, and we present it as such rather than as a benchmarked
result. The per-image binary decision is the maximum-confidence pain-box score
compared against a tuned threshold (Section~\ref{sec:eng-detector-imbalance});
the DETR-internal transformer backbone is, like the rest of the engine,
conceded plumbing.

\subsection{FGS-safe augmentation}
\label{sec:eng-detector-aug}

The subtle orbital, muzzle, and whisker cues that the Feline Grimace Scale
reads are the entire signal, and any augmentation that distorts facial geometry
poisons both the detector and the crop it produces. We therefore restrict
augmentation to label- and geometry-preserving transforms. Horizontal flip is
permitted (the face is left/right symmetric); small translation and scale and
rotations up to $\sim 10^{\circ}$ are permitted; modest saturation/value jitter
is permitted. Vertical flip is banned outright, because it inverts grimace
geometry; large rotations, heavy colour distortion, and aggressive
mixup/copy-paste \emph{blends} are banned for the same reason. Hue jitter is
held near zero, since coat colour is a capture-condition confound we do not wish
to amplify (the confound itself is audited in a later chapter). Mosaic is
enabled early and disabled for the final epochs so that the last training views
match the real distribution.

\subsection{Imbalance ladder and operating point}
\label{sec:eng-detector-imbalance}

The corpus is heavily imbalanced toward the no-pain class. On the metadata
basis we report throughout --- $264$ pain boxes against $1{,}819$ no-pain boxes,
a prevalence of $\approx 12.7\%$ and an imbalance of $\approx 6.9{:}1$ --- we
note explicitly that these are \emph{box counts, not individual cats}, and that
the live-annotation basis ($246/1{,}811$, $\approx 12.0\%$) is logged as a
discrepancy rather than silently reconciled. We address imbalance with data and
loss before architecture, climbing a fixed ladder and changing one lever per
run, with the decision threshold tuned \emph{last}
(Table~\ref{tab:eng-imbalance}).

\begin{table}[t]
\centering
\footnotesize
\caption[Imbalance ladder for the Phase~A detector]{Imbalance ladder, climbed
cheapest-first; one lever changed per run; the decision threshold moved only
after the data/loss rungs are exhausted. The $\mathrm{pos\_weight}\approx 2.6$
basis is the metadata box ratio $\sqrt{1819/264}$ (box counts, not
individuals); the live-annotation alternative ($\approx 2.7$) is acceptable when
labelled by basis. Values are design defaults, not measured results.}
\label{tab:eng-imbalance}
\begin{tabular}{@{}clll@{}}
\toprule
Rung & Lever & Setting & Fires when \\
\midrule
1 & Class weight / $\mathrm{pos\_weight}$
  & $\sqrt{1819/264}\approx 2.6$ (cap $\sim 3$)
  & always (baseline) \\
2 & Oversample + object copy-paste
  & duplicate pain images; paste pain crops
  & rung-1 recall below target \\
3 & Focal loss
  & $\gamma=1.5,\ \alpha=0.25$
  & rungs 1--2 insufficient \\
4 & Threshold tuning (last)
  & lowest threshold meeting recall $\geq 0.90$
  & only after 1--3 \\
\bottomrule
\end{tabular}
\end{table}

The operating point is fixed at \textbf{pain recall $\geq 0.90$}, selected
inside training folds on the validation precision--recall curve --- explicitly
\emph{not} a Youden-$J$ or $F_1$ knee. A symmetric-cost operating point is the
wrong loss for a welfare instrument in which undertreatment is far costlier than
overtreatment; the welfare-asymmetric decision curve that exploits this choice
is the wrapper's headline artifact and is developed later. The $\geq 0.90$ floor
is justified against the $90.7\%$ reference sensitivity of the Feline Grimace
Scale validation \citep{evangelista2019}, not against prior automated work, and
the specificity it buys is reported with bootstrap confidence intervals across
the five cat-grouped folds. Throughout, augmented and oversampled copies never
enter any reported $N$; every rate is accompanied by the distinct-pain-cat
denominator and Clopper--Pearson or bootstrap intervals, and accuracy and
ROC-AUC are never used as the ranking metric. The primary detection metric is
the precision--recall AUC (average precision) with pain as the positive class.

\subsection{The circularity firewall at the operating point}
\label{sec:eng-detector-firewall}

Phase~A is where the project's central methodological guard first bites. The
sensitivity and specificity reported at the operating point are estimated
\emph{only} on vet-confirmed binary labels. Agreement of any model output
against the VLM-derived weak labels is never treated as validation, because it
would measure the model re-learning the VLM's heuristic rather than tracking
ground truth. This firewall is implemented as a structural constraint on which
label column the evaluation code may read, and it is the reason the binary
decision --- and not the graded sum --- carries the validated claim.

\section{Face-crop and eye-alignment preprocessing (Gate~5)}
\label{sec:eng-crop}

The detector box is a face-crop \emph{candidate}, not a validated cropper. This
section specifies the detect $\rightarrow$ quality-gate $\rightarrow$ expand
$\rightarrow$ align $\rightarrow$ resize pipeline that turns a candidate box
into the fixed crop tensor the backbone consumes, and the blocking Gate~5 audit
that decides whether the cheap box crop suffices or whether a two-point
eye-similarity aligner must be switched on. The motivation is defensive: a
misaligned crop makes any downstream ``calibration'' number a measure of crop
quality rather than of pain, so this stage exists to stop us mistaking one for
the other.

\subsection{Margin expansion}
\label{sec:eng-crop-expand}

The ear AU is read from ear tips and the whisker AU from whisker splay, both of
which sit \emph{outside} a tight head box. We therefore \emph{expand} the box
before cropping rather than shrink it: a $\sim 13.5\%$ margin (midpoint of a
$12$--$15\%$ band), applied asymmetrically with extra room above the head to
recover ear tips. A box already abutting the frame edge cannot be expanded to
recover a clipped ear or whisker; such cases are flagged and routed to the vet
queue as an abstention signal rather than scored on a silently-truncated AU.

\subsection{Frontal-pose and quality gate}
\label{sec:eng-crop-quality}

The Feline Grimace Scale assumes a near-frontal view. Profile, tilted,
occluded, and eyes-closed faces are not scored; they are routed to the vet
queue and never enter the feature cache or any reported denominator. The gate is
a deterministic, training-free, CPU-only filter over the expanded crop: a
yaw/frontality check from the eye-pair geometry, an in-plane tilt check from the
inter-ocular line (the aligner can correct up to $\sim 25^{\circ}$ of roll;
beyond that is a pose problem, not an alignment problem), an eye-aspect-ratio
check for closed eyes, a variance-of-Laplacian blur check, and a coarse Haar
\texttt{frontalcatface} detector whose failure to find a face is itself an
abstention signal. The Haar cascade is used only as a zero-training fallback
localiser and abstention signal --- never as a landmark pipeline, an explicit
avoidance of the landmark-then-geometry trap of prior automated work
\citep{steagall2023}. Each decision and its reasons are written to the crop
manifest; all thresholds are tuned once on the audit set and then frozen, so the
gate cannot be re-tuned to inflate apparent prevalence.

\subsection{The Gate~5 NME audit and the eye-alignment fallback}
\label{sec:eng-crop-gate5}

Gate~5 is blocking: no ordinal head is trained and no vet hour is spent until it
passes. On a small audit set ($30$--$50$ images) it answers two yes/no
questions, using the $48$ CatFLW facial landmarks \citep{martvel2024catflw} as
ground truth. First, is the interocular-normalised mean landmark error (NME)
inside the box crop low enough, and is the box-vs-aligned NME gap small enough,
that the box crop can feed the backbone directly? Second, is the median face
resolution after resize to the backbone's patch grid high enough that AU detail
survives? Design thresholds are a median NME $\leq 0.08$ with a box-vs-aligned
gap $\leq 0.02$, and a median interocular distance $\geq 40$\,px after resize;
these are proposed operating values calibrated to the automated-landmark NME
range of $9$--$26\%$ reported across morphologies by \citet{martvel2024catflw},
to be confirmed against the audit output before freezing. The same work reports
a roughly seven-point pain-accuracy cost from automated versus manual landmarks,
which is precisely the failure mode this gate is built to detect.

If the audit fails, a closed-form \emph{two-point eye-similarity} transform is
switched on as the production crop step: a rotate-plus-uniform-scale-plus-%
translate map taking the two detected eye centres onto fixed canonical targets.
It is a similarity transform with no learned warp and no per-image optimisation,
computed on CPU. Whichever transform Gate~5 certifies --- raw box crop or
eye-aligned crop --- is then frozen and applied once to the full corpus,
producing $518\times518$ RGB crops (with $518 = 14\times 37$ divisible by the
backbone patch size). Crops flagged for the vet queue are excluded from the
feature cache and from every sensitivity/specificity/$\kappa$ denominator,
turning ``detector or quality failure $\rightarrow$ defer to vet'' into an
explicit, counted abstention channel rather than a silent drop.

\section{The frozen DINOv2 ViT-S/14 backbone and feature caching}
\label{sec:eng-backbone}

The encoder is a frozen DINOv2 ViT-S/14 \citep{oquab2023dinov2}: hidden width
$384$, patch size $14$, forward-only. The \emph{register} variant
\texttt{dinov2\_vits14\_reg} is the field default, because registers suppress the
attention artifacts that hurt dense, localized features --- exactly the per-AU
regime (orbital, ear, muzzle); the reg variant is A/B'd against the plain variant
(kept for backwards compatibility) before ViT-S/14 is locked. We freeze it because the supervised set is
on the order of $120$--$300$ vet- and VLM-labelled faces, a regime in which
full fine-tuning of a vision transformer would be malpractice; only the light
ordinal heads are trained. Patch size and hidden width are derived from the
loaded model rather than hard-coded in load-bearing code. The default pooled
representation for the heads is the class token; the mean of the patch tokens is
cached as a second column so that the class-token-versus-mean-patch comparison
is a configuration flag rather than a re-cache.

The single most important engineering decision in this section is
\emph{feature caching}, and it is the decision that makes the whole engine
tractable on an Apple~M4 / MPS machine with no CUDA. Because the backbone is
frozen, its output for a given crop never changes; we therefore run exactly one
forward pass over every crop, write the pooled features to disk, and train the
heads off the cache with \emph{zero} backbone forwards per epoch. This turns
each head epoch from minutes into milliseconds and makes the head sweep
(pooling $\times$ loss $\times$ seeds) a matter of minutes rather than hours.
The DINOv2 commit hash, the preprocessing parameters, the device, and the torch
version are recorded alongside the cache, because the cache is an input to every
reported number. We do \emph{not} domain-adapt the backbone on any external
archive: at this label budget the risk of degrading the features for near-zero
gain is real, and the backbone's generality is what we are relying on.

\section{The five per-AU CORN ordinal heads}
\label{sec:eng-corn}

Each of the five FGS action units --- ear, orbital tightening, muzzle, whiskers,
and head position --- is an ordinal variable on the levels $\{0,1,2\}$. We model
each with a \emph{rank-consistent} ordinal regression head trained with the CORN
(Conditional Ordinal Regression for Neural networks) loss
\citep{shi2023corn}. CORN is chosen over a plain softmax classifier because the
levels are ordered, and over CORAL \citep{cao2020coral} because CORN attains
rank consistency without CORAL's shared-bias capacity restriction. With three
levels per AU, each head emits $K-1 = 2$ conditional logits, so the five AUs
together require a single $384 \rightarrow 10$ linear projection whose output is
split into five two-logit chunks --- one matmul, then a \texttt{split}. The head
is a dropout-regularised linear probe with small-variance truncated-normal
initialisation appropriate to the tiny-data linear-probe regime. The CORN loss
itself is vendored as roughly forty auditable lines (torch and functional only,
no pip dependency) so that the datasheet shows exactly what trained the heads;
it is cross-checked once against a reference implementation in the Gate~4 unit
test (Section~\ref{sec:eng-decode-gate4}) and the reference dependency is then
dropped.

A sibling binary pain head --- a $384 \rightarrow 1$ linear layer trained with a
weighted binary cross-entropy on the pain label --- shares the same cache, is
independent of the CORN heads, and is the v1 spine that the trustworthiness
wrapper actually operates on. The five-AU sum is the inspected path.

The decode contract that turns logits into a level, a sum, and a decision is the
same in this chapter, in the VLM weak-labelling pipeline, and in the wrapper:
the hard level is $\sum_k \mathbb{1}[\operatorname{cumprod}(\sigma(\text{logits}))_k > 0.5]$,
the sum is the sum of the five AU levels in $\{0,\dots,10\}$, and the decision
fires when the ratio $\text{sum}/10$ reaches the clinical cut of $0.39$
(approximately $4/10$) reported by the FGS validation \citep{evangelista2019}.
Crucially the \emph{VLM never emits the sum or the decision}: it emits the five
ordinal atoms, and the sum and the $0.39$ flag are computed in code, so that
there is a single threshold definition in the repository. Figure~\ref{fig:eng-corn}
shows the head and its two decode paths.

\begin{figure}[t]
\centering
\begin{tikzpicture}[
  node distance=5mm and 8mm,
  box/.style={draw, rounded corners, align=center, minimum height=8mm,
              inner sep=3pt, font=\footnotesize, text width=22mm},
  soft/.style={box, fill=black!4},
  hard/.style={box, draw, very thick},
  arr/.style={-{Latex[length=2mm]}, thick},
  lbl/.style={font=\scriptsize\itshape, text=black!55}
]
\node[box] (feat) {cached feature\\$\mathbf{z}\in\mathbb{R}^{384}$};
\node[box, right=of feat] (proj) {linear\\$384\!\to\!10$};
\node[box, right=of proj] (split) {split\\$5\times[\,\cdot\,,2]$};
\node[box, right=of split] (cum) {$\operatorname{cumprod}\sigma$\\$=P(\text{rank}>k)$};

\node[soft, below=10mm of cum] (pmf) {per-AU pmf\\$(p_0,p_1,p_2)$};
\node[soft, left=14mm of pmf] (conv) {convolve 5\\$\to$ 0--10 pmf};
\node[soft, left=of conv] (rps) {RPS-on-sum\\+ ClasswiseECE\\\emph{(inspect)}};

\node[hard, above=10mm of cum] (thr) {$>0.5 \Rightarrow$\\level $\in\{0,1,2\}$};
\node[hard, left=14mm of thr] (sum) {sum 0--10\\$\geq 0.39$};

\draw[arr] (feat) -- (proj);
\draw[arr] (proj) -- (split);
\draw[arr] (split) -- (cum);
\draw[arr] (cum) -- (pmf);
\draw[arr] (pmf) -- (conv);
\draw[arr] (conv) -- (rps);
\draw[arr] (cum) -- (thr);
\draw[arr] (thr) -- (sum);

\node[lbl, above=0.5mm of thr.north] {point-decision path (decision only)};
\node[lbl, below=0.5mm of rps.south, text width=30mm, align=center]
  {distributional path (default; inspected-not-validated)};
\end{tikzpicture}
\caption[CORN head and distributional decode]{The rank-consistent CORN head and
its two decode paths. A cached feature is projected to five two-logit chunks;
the cumulative product of the sigmoids gives, per AU, $P(\text{rank}>k)$. The
\emph{soft} path (lower, default) turns these into a per-AU pmf and convolves
the five into a distribution over the $0$--$10$ sum, scored by the ranked
probability score and per-AU classwise ECE; this path is inspected, not
validated. The \emph{hard} path (upper) thresholds the cumulative products at
$0.5$ to a level, sums to $0$--$10$, and fires the decision at the $0.39$ ratio.
The hard path is used for the point decision only.}
\label{fig:eng-corn}
\end{figure}

\section{Distributional decode}
\label{sec:eng-decode}

The default scoring path is distributional, not argmax. From each AU's soft
cumulative probabilities $P(\text{rank}>k) = \operatorname{cumprod}(\sigma(\text{logits}))$
we form a proper probability mass function over $\{0,1,2\}$ via
$p_0 = 1 - P(\text{rank}>0)$, $p_1 = P(\text{rank}>0) - P(\text{rank}>1)$, and
$p_2 = P(\text{rank}>1)$, clamping tiny negative values from floating-point
error. The five per-AU pmfs are then \emph{convolved} into a single pmf over the
$0$--$10$ sum. From that distribution we report two things: the ranked
probability score on the sum, as a single scalar with a bootstrap confidence
interval, and a per-AU classwise expected calibration error
\citep{guo2017calibration}. The argmax / hard decode is reserved for the $0.39$
point decision alone.

One reporting decision in this path is load-bearing and we state it explicitly:
\emph{we do not draw a binned reliability diagram on the eleven-atom $0$--$10$
sum}. With only eleven atoms and the label budget we have, per-bin standard
errors are large enough ($\pm 0.18$--$0.26$) that such a diagram would launder
noise as evidence. The ranked probability score on the sum plus per-AU classwise
ECE is the statistically coherent substitute, and it is reported as
inspected-not-validated evidence about the graded layer, never as a validated
calibration claim.

\subsection{Gate~4: blocking decode correctness}
\label{sec:eng-decode-gate4}

Before any score from the engine is believed, two correctness checks must pass,
and they are blocking in continuous integration. The first is a synthetic
decode unit test: hand-verified logit configurations are asserted to decode to
known levels, the decode is asserted monotone in the logits (raising any logit
cannot lower the decoded level, the rank-consistency property), the per-AU pmf
is asserted to sum to one with no negative atoms, the convolved sum pmf is
asserted to be a distribution over $0$--$10$, and the $0.39$ point decision is
asserted at its boundary (a sum of $4$ fires, a sum of $3$ does not). The second
is an MPS-versus-CPU logit-parity check on a real aligned crop, guarding against
silently-wrong accelerator arithmetic --- a leak-proof metric computed from
wrong logits is worthless. A red Gate~4 blocks the entire engine; the vendored
decode is additionally cross-checked once against a reference CORN implementation,
and on disagreement the vendoring is fixed rather than the expected value
adjusted to match a broken decode.

\section{Training: co-teaching small-loss on VLM mass plus a clean vet anchor}
\label{sec:eng-training}

The heads are trained on a mixture of two label sources of very different
quality: a large mass of VLM weak labels covering all scored crops, and a small
clean anchor of vet-confirmed labels (on the order of $120$--$300$ faces, at
least $50$ pain-positive, with the exact integer fixed by the project's power
calculation before any vet hour is spent). Training on the clean anchor alone
gives too few examples; training naively on all labels bakes in the VLM's
under-estimation bias. We therefore use \emph{co-teaching with small-loss
selection}: two heads are trained in parallel, each selecting the small-loss
subset of each mini-batch for its partner, following the canonical Co-teaching
schedule $R(T)=1-\tau\cdot\min(T/T_k,1)$~\cite{han2018coteaching}: the keep-rate
ramps from inclusive toward selective over the first $T_k$ epochs, starting at
$1.0$ (keep all) and descending to $1-\tau$, where $\tau$ is the estimated label
noise rate --- exploiting the memorization effect, whereby networks fit clean
patterns before noisy ones, so noisy examples are pruned late rather than early.
Vet-clean rows are never
dropped from either head's selection --- they are protected anchor points that
the noisy mass is filtered against. AUs left unscored on a given image carry a
sentinel and are skipped in that image's loss, and the weakest AUs (muzzle and
whiskers) may optionally be upweighted. Figure~\ref{fig:eng-coteach} sketches the
scheme.

\begin{figure}[t]
\centering
\fbox{\begin{minipage}{0.92\linewidth}
\centering\footnotesize
\vspace{2mm}
\begin{tikzpicture}[
  node distance=5mm and 10mm,
  blk/.style={draw, rounded corners, align=center, font=\footnotesize,
              minimum height=8mm, inner sep=3pt, text width=24mm},
  arr/.style={-{Latex[length=2mm]}, thick},
  dasharr/.style={-{Latex[length=2mm]}, thick, dashed}
]
\node[blk] (mass) {VLM weak mass\\(noisy, all crops)};
\node[blk, below=4mm of mass] (anchor) {vet-clean anchor\\(protected)};
\node[blk, right=20mm of mass] (A) {head A};
\node[blk, right=20mm of anchor] (B) {head B};
\draw[arr] (mass) -- (A);
\draw[arr] (mass) -- (B);
\draw[arr] (anchor.east) -- ($(A.west)+(0,-2mm)$);
\draw[arr] (anchor.east) -- ($(B.west)+(0,2mm)$);
\draw[dasharr] (A.south) to[bend left=18] node[midway, right, font=\scriptsize\itshape]
   {A's small-loss subset trains B} (B.north);
\draw[dasharr] (B.north) to[bend left=18] node[midway, left, font=\scriptsize\itshape]
   {B's small-loss subset trains A} (A.south);
\end{tikzpicture}
\vspace{1mm}

\emph{Placeholder schematic --- to be replaced by the produced training-loss /
keep-rate figure once runs exist.}
\vspace{2mm}
\end{minipage}}
\caption[Co-teaching small-loss training]{Co-teaching small-loss training. Two
heads consume the same noisy VLM mass; each selects the small-loss subset of a
mini-batch to train its partner, so each head is shielded from the other's
memorised label noise. The vet-clean anchor rows are protected and never
filtered out. Captioned placeholder; no loss curve is shown because no training
run has been performed.}
\label{fig:eng-coteach}
\end{figure}

Folds are cat-disjoint, read from the per-individual merge and the frozen,
hashed hold-out manifest, so that no individual straddles a fold and augmented
copies never enter a reported $N$. The head sweep --- class-token versus
mean-patch pooling, CORN versus CORAL loss, five seeds with mean and bootstrap
CI --- runs in minutes precisely because the backbone never re-runs.

\section{What the engine outputs, and what is and is not claimed}
\label{sec:eng-outputs}

The engine emits the artifacts in Table~\ref{tab:eng-outputs}: cached features,
trained CORN heads, a trained binary pain head, per-AU soft cumulative
probabilities on the hold-out, a convolved $0$--$10$ sum pmf, and the blocking
Gate~4 report. The claim discipline attached to each is the point of this
chapter. The binary pain decision plus abstention is the validated spine, with
sensitivity and specificity estimated downstream on vet-confirmed labels only.
The $0$--$10$ sum is inspected-not-validated: its agreement against the VLM is
never reported as validation, it is discussed qualitatively with the circularity
stated, and the project's abstract holds with it dropped. The engine is never
claimed novel --- a reviewer applying the ``swap the backbone, add some
metrics'' filing finds nothing here to kill, by design. And the kill/pivot rule
is wired in: if the per-AU agreement of the orbital, ear, and head heads against
the vet anchor (gated in a later chapter on the confidence-interval lower bound,
not the point estimate) falls below its pre-registered floor, the graded layer
is dropped entirely and the system ships as a calibrated binary detector with
abstention and the confound audit; the engine's heads still train and ship as
the inspected artifact, but no graded claim is made.

\begin{table}[t]
\centering
\footnotesize
\caption[Phase~B engine outputs and claim status]{Engine outputs and their claim
status. The binary spine is the only validated output; the graded sum and its
pmf are inspected-not-validated; everything else is intermediate plumbing. No
output in this table carries a measured result in this document.}
\label{tab:eng-outputs}
\begin{tabular}{@{}lll@{}}
\toprule
Artifact & Role & Claim status \\
\midrule
Cached frozen features        & input to all heads      & input only \\
Trained CORN heads (5 seeds)  & per-AU ordinal scores   & conceded plumbing \\
Binary pain head              & v1 spine                & \textbf{validated} (vet labels only) \\
Per-AU soft cumulative probs  & feeds wrapper           & intermediate \\
Convolved $0$--$10$ sum pmf   & graded layer            & \emph{inspected-not-validated} \\
Gate~4 decode/parity report   & correctness            & blocking pass/fail \\
\bottomrule
\end{tabular}
\end{table}

In sum, this chapter has specified a deliberately unremarkable engine ---
detector, alignment-gated crop, frozen backbone with cached features,
rank-consistent ordinal heads, distributional decode, and noise-robust training
--- whose every design choice is the conservative one for a low-data,
single-machine, single-vet regime. Its outputs feed the headline confound-attribution
protocol, the guarded VLM-as-AU-rater reliability check, and the cited
trustworthiness wrapper that carry this work's contribution. Nothing here
is offered as novel, and nothing here is reported as a result.

The portable method surface --- the headline capture-condition
confound-attribution protocol and the supporting VLM-as-AU-rater $\kappa$
reliability check --- is dataset-agnostic code in \texttt{src/protocols/}
(with thin adapters for AU override, column mapping, and generic non-FGS ordinal
mode; see \texttt{src/protocols/adapters.py} and re-exports in
\texttt{src/protocols/\_\_init\_\_.py}). The standalone test corpus
(\texttt{src/protocols/standalone\_test\_corpus.py}) exercises the surface on
synthetic arbitrary-AU data with zero torch leakage and is deletion-safe
(imports only the protocols layer); the headline confound leg is additionally
validated today on planted positive and negative controls. The engine remains
conceded plumbing. See also the evaluation protocol
chapter for the gate-enforced usage contract and \texttt{Makefile} targets
\texttt{test-portable} / \texttt{gate-e2e-synthetic}.

The protocols seam (portable pmf derives re-exported del-safe from decode for arbitrary AU via adapters/generic_ordinal_mode; N/k from len/shape, no hard 5/11 in portable path), generalized decode (engine path), DINOv3 prep, and cache schema are the
executable embodiments of portable: use \texttt{src.protocols} adapters +
\texttt{get\_au\_names(override)} + \texttt{generic\_ordinal\_mode} (no 0.39 assumption) for arbitrary ordinal corpora; point-decision
0.39 remains FGS-compat only (single-source in constants). The frozen backbone and feature cache
(\texttt{src/model/backbone.py}, \texttt{src/model/cache\_features.py}) are
DINOv3-ready (SUPPORTED\_VARIANTS includes dinov3\_*; richer patch\_std pooling
per fresh MCP context7 /facebookresearch/dinov3 for localized AU cues; explicit
\_CACHE\_SCHEMA + FEATURE\_SCHEMA with version/n\_aus/k/feature\_dim/patch\_mode
+ full PROVENANCE.json for portable repro across backbones; dinov commit/device/torch
captured). These, with the central orchestrator and standalone exerciser, make the
``portable method'' claim accurate and citable. The engine itself stays
conceded plumbing.
